\documentclass{report}
\usepackage{amsmath,amssymb}
\setlength{\parindent}{0mm}
\setlength{\parskip}{1em}
\begin{document}
\begin{center}
\rule{6in}{1pt} \
{\large Jed Brown \\
{\bf HPGMG: Benchmarking computers using multigrid}}

Argonne National Lab \& CU Boulder
\\
{\tt jed@jedbrown.org}\\
Mark Adams\\
Sam Williams\end{center}

HPGMG (https://hpgmg.org) is a geometric multigrid benchmark designed to
measure the performance and versatility of computers. For a benchmark to
be representative of applications, good performance on the benchmark
should be sufficient to ensure good performance on most important
applications and only those system features necessary for some important
applications should be stressed by the benchmark. Moreover, the
specification should be scale-free and the method should solve a
globally-coupled problem. FMG with highly parallel smoothers and a
matrix-free representation of the operator has most of the desirable
attributes and we have developed benchmark implementations using finite
volume and finite element methods. Both implementations stress the
network similarly, and can clearly distinguish different network types,
but have quite different on-node characteristics due to different
arithmetic intensity, cache demands, and the fact that a non-overlapping
partition of the data is not a partition of the work. In this talk, I
will present the benchmark design, performance results on many top
machines, and remaining challenges for this benchmarking effort.


\end{document}
